\subsection{Representation Learning Objectives}

Representation learning trains maps that make downstream prediction, reconstruction, comparison, or generation easier.

\begin{definition} \label{def:sl-autoencoder} \hlt{Autoencoder}\\
An autoencoder is an encoder-decoder pair $(e,d)$ trained to minimize a reconstruction criterion such as
\[
\frac1n\sum_{i=1}^n\norm{x_i-d(e(x_i))}_2^2.
\]
Bottlenecks, noise, sparsity, or probabilistic latent variables prevent the identity map from being the only useful solution.
\end{definition}

\begin{definition} \label{def:sl-sparse-representation-objective} \hlt{Sparse Representation Objective}\\
Given a dictionary $W$, sparse coding seeks representations $h_i$ with small reconstruction error and many zero coordinates:
\[
\min_{W,h_1,\ldots,h_n}
\sum_{i=1}^n\norm{x_i-Wh_i}_2^2+\lambda\norm{h_i}_1.
\]
\end{definition}

\begin{definition} \label{def:sl-contrastive-objective} \hlt{Contrastive Objective}\\
Contrastive learning trains representations so related examples have high similarity while unrelated examples have low similarity.  The negative examples and augmentation distribution are part of the statistical model of the task.
\end{definition}

\begin{remark} \label{rem:sl-representation-source-migration} \hlt{From Regularization Source to Representation Learning}\\
The old sparse-representation notes appeared inside deep-learning regularization.  They belong here as representation objectives, while their neural-network penalties reappear in the neural-network regularization section.
\end{remark}
